% GENERATED FILE. Edit canonical lessons, then run npm run generate:books.
% canonical-source-hash: fnv1a64:36f1c1401288e9ce
% canonical-lessons: BN-C297-kotha, BN-C297-sakkhatkar, BN-C297-ghoshona-kora, BN-C297-bibriti, BN-C297-ullekh-kora

\chapter{Who Said What}
\label{ch:bn-who-said-what}
\hlchaptermodality{297}

\begin{chapteropening}
\textbf{By the end of this chapter:} \emph{I can report what someone said, announced or mentioned.}

Complete the last lesson of chapter 297: I can report what someone said, announced or mentioned.
\end{chapteropening}

\section[\texorpdfstring{\bn{কথা} (kôthā)}{কথা (kôthā)}]{\bn{কথা} (kôthā) — words; talk; what someone said}

\label{lesson:BN-C297-kotha}



\begin{quote}
Before the new one: say the Bengali for to publish, then the Bengali for an editor.
\end{quote}

\subsection*{You'll want to know: \bn{কথা}}
\textbf{\bn{কথা}} — \emph{kôthā} — \textquotedblleft{}words; talk; what someone said\textquotedblright{}.

\begin{quote}
\bn{সাক্ষীর} \bn{কথা} \bn{শুনে} \bn{দারোগা} \bn{কিছু} \bn{বললেন} \bn{না।}
\end{quote}

\begin{quote}
\emph{sākkhīr kôthā shune dārogā kichhu bôllen nā.}
\end{quote}

\begin{quote}
The inspector listened to the witness and said nothing.
\end{quote}

\bn{কথা} \bn{বলা}, to talk; \bn{কথা} \bn{দেওয়া}, to give your word.

\begin{cousinweb}
\textbf{\bn{কথা}} (\emph{kôthā}) is Sanskrit \emph{kathā}, 'talk, conversation; a story' (Monier-Williams), and it keeps both senses in Bengali.
\end{cousinweb}

\begin{grammarlens}[title={What you've built}]
One more word for telling what happened.
\end{grammarlens}

\subsection*{Your turn}
\begin{itemize}
\raggedright
  \item \emph{Say it:} \emph{kôthā}
  \item \emph{Say it:} \emph{kôthā}, once more
  \item \emph{Say it:} \emph{sākkhīr kôthā shune dārogā kichhu bôllen nā.} — The inspector listened to the witness and said nothing.
  \item \emph{Recall:} say the Bengali for to publish, then the Bengali for an editor, then say \emph{kôthā} again
\end{itemize}

\begin{tcolorbox}[breakable,colback=teal!4,colframe=teal!35!black,title={Before you move on}]
What does \bn{কথা} mean? (\textquotedblleft{}Words; talk; what someone said\textquotedblright{}.) Say it once more.
\end{tcolorbox}
\section[\texorpdfstring{\bn{সাক্ষাৎকার} (sākkhātkār)}{সাক্ষাৎকার (sākkhātkār)}]{\bn{সাক্ষাৎকার} (sākkhātkār) — an interview}

\label{lesson:BN-C297-sakkhatkar}



\begin{quote}
Before the new one: say the Bengali for an editor, then the Bengali for talk.
\end{quote}

\subsection*{You'll want to know: \bn{সাক্ষাৎকার}}
\textbf{\bn{সাক্ষাৎকার}} — \emph{sākkhātkār} — \textquotedblleft{}an interview\textquotedblright{}.

\begin{quote}
\bn{টিভিতে} \bn{সাক্ষীর} \bn{একটা} \bn{সাক্ষাৎকার} \bn{দেখাল।}
\end{quote}

\begin{quote}
\emph{ṭibhite sākkhīr ekṭā sākkhātkār dekhālô.}
\end{quote}

\begin{quote}
The television showed an interview with the witness.
\end{quote}

\begin{cousinweb}
\textbf{\bn{সাক্ষাৎকার}} (\emph{sākkhātkār}) is Sanskrit \emph{sākṣātkāra}, 'a direct perceiving' (Monier-Williams), from \emph{sākṣāt}, 'before one's eyes'; the sense 'an interview', a meeting face to face, is the modern one.
\end{cousinweb}

\begin{grammarlens}[title={What you've built}]
One more word for telling what happened.
\end{grammarlens}

\subsection*{Your turn}
\begin{itemize}
\raggedright
  \item \emph{Say it:} \emph{sākkhātkār}
  \item \emph{Say it:} \emph{sākkhātkār}, once more
  \item \emph{Say it:} \emph{ṭibhite sākkhīr ekṭā sākkhātkār dekhālô.} — The television showed an interview with the witness.
  \item \emph{Recall:} say the Bengali for an editor, then the Bengali for talk, then say \emph{sākkhātkār} again
\end{itemize}

\begin{tcolorbox}[breakable,colback=teal!4,colframe=teal!35!black,title={Before you move on}]
What does \bn{সাক্ষাৎকার} mean? (\textquotedblleft{}An interview\textquotedblright{}.) Say it once more.
\end{tcolorbox}
\section[\texorpdfstring{\bn{ঘোষণা} \bn{করা} (ghoshônā kôrā)}{ঘোষণা করা (ghoshônā kôrā)}]{\bn{ঘোষণা} \bn{করা} (ghoshônā kôrā) — to announce; to declare}

\label{lesson:BN-C297-ghoshona-kora}



\begin{quote}
Before the new one: say the Bengali for talk, then the Bengali for an interview.
\end{quote}

\subsection*{You'll want to know: \bn{ঘোষণা} \bn{করা}}
\textbf{\bn{ঘোষণা} \bn{করা}} — \emph{ghoshônā kôrā} — \textquotedblleft{}to announce; to declare\textquotedblright{}.

\begin{quote}
\bn{সরকার} \bn{দুই} \bn{দিনের} \bn{ছুটি} \bn{ঘোষণা} \bn{করল।}
\end{quote}

\begin{quote}
\emph{sôrkār dui diner chhuṭi ghoshônā kôrlô.}
\end{quote}

\begin{quote}
The government announced a two-day holiday.
\end{quote}

Past (he/she …): \bn{তিনি} \bn{ফলাফল} \bn{ঘোষণা} \bn{করলেন}, \emph{tini phôlāphôl ghoshônā kôrlen}.

\begin{cousinweb}
\textbf{\bn{ঘোষণা} \bn{করা}} (\emph{ghoshônā kôrā}) puts Sanskrit \emph{ghoṣaṇā}, 'a proclaiming aloud' (Monier-Williams), before \bn{করা}: to make a thing known to everyone at once.
\end{cousinweb}

\begin{grammarlens}[title={What you've built}]
One more word for telling what happened.
\end{grammarlens}

\subsection*{Your turn}
\begin{itemize}
\raggedright
  \item \emph{Say it:} \emph{ghoshônā kôrā}
  \item \emph{Say it:} \emph{ghoshônā kôrā}, once more
  \item \emph{Say it:} \emph{sôrkār dui diner chhuṭi ghoshônā kôrlô.} — The government announced a two-day holiday.
  \item \emph{Recall:} say the Bengali for talk, then the Bengali for an interview, then say \emph{ghoshônā kôrā} again
\end{itemize}

\begin{tcolorbox}[breakable,colback=teal!4,colframe=teal!35!black,title={Before you move on}]
What does \bn{ঘোষণা} \bn{করা} mean? (\textquotedblleft{}To announce; to declare\textquotedblright{}.) Say it once more.
\end{tcolorbox}
\section[\texorpdfstring{\bn{বিবৃতি} (bibriti)}{বিবৃতি (bibriti)}]{\bn{বিবৃতি} (bibriti) — a public statement; a statement}

\label{lesson:BN-C297-bibriti}



\begin{quote}
Before the new one: say the Bengali for an interview, then the Bengali for to announce.
\end{quote}

\subsection*{You'll want to know: \bn{বিবৃতি}}
\textbf{\bn{বিবৃতি}} — \emph{bibriti} — \textquotedblleft{}a public statement; a statement\textquotedblright{}.

\begin{quote}
\bn{মন্ত্রী} \bn{দুর্ঘটনার} \bn{পরে} \bn{একটা} \bn{বিবৃতি} \bn{দিলেন।}
\end{quote}

\begin{quote}
\emph{môntrī durghôṭônār pôre ekṭā bibriti dilen.}
\end{quote}

\begin{quote}
The minister made a statement after the accident.
\end{quote}

\begin{cousinweb}
\textbf{\bn{বিবৃতি}} (\emph{bibriti}) is Sanskrit \emph{vivṛti}, 'a making manifest, an exposition' (Monier-Williams): a statement opens a matter up. \bn{বিবৃতি} \bn{দেওয়া}, to issue a statement.
\end{cousinweb}

\begin{grammarlens}[title={What you've built}]
One more word for telling what happened.
\end{grammarlens}

\subsection*{Your turn}
\begin{itemize}
\raggedright
  \item \emph{Say it:} \emph{bibriti}
  \item \emph{Say it:} \emph{bibriti}, once more
  \item \emph{Say it:} \emph{môntrī durghôṭônār pôre ekṭā bibriti dilen.} — The minister made a statement after the accident.
  \item \emph{Recall:} say the Bengali for an interview, then the Bengali for to announce, then say \emph{bibriti} again
\end{itemize}

\begin{tcolorbox}[breakable,colback=teal!4,colframe=teal!35!black,title={Before you move on}]
What does \bn{বিবৃতি} mean? (\textquotedblleft{}A public statement; a statement\textquotedblright{}.) Say it once more.
\end{tcolorbox}
\section[\texorpdfstring{\bn{উল্লেখ} \bn{করা} (ullekh kôrā)}{উল্লেখ করা (ullekh kôrā)}]{\bn{উল্লেখ} \bn{করা} (ullekh kôrā) — to mention}

\label{lesson:BN-C297-ullekh-kora}



\begin{quote}
Before the new one: say the Bengali for to announce, then the Bengali for a public statement.
\end{quote}

\subsection*{You'll want to know: \bn{উল্লেখ} \bn{করা}}
\textbf{\bn{উল্লেখ} \bn{করা}} — \emph{ullekh kôrā} — \textquotedblleft{}to mention\textquotedblright{}.

\begin{quote}
\bn{প্রতিবেদনে} \bn{সাংবাদিক} \bn{সাক্ষীর} \bn{নাম} \bn{উল্লেখ} \bn{করলেন} \bn{না।}
\end{quote}

\begin{quote}
\emph{prôtibedône sāngbādik sākkhīr nām ullekh kôrlen nā.}
\end{quote}

\begin{quote}
In the report the journalist did not mention the witness's name.
\end{quote}

Past (he/she …): \bn{তিনি} \bn{ঘটনাটা} \bn{উল্লেখ} \bn{করলেন}, \emph{tini ghôṭônāṭā ullekh kôrlen}.

\begin{cousinweb}
\textbf{\bn{উল্লেখ} \bn{করা}} (\emph{ullekh kôrā}) puts Sanskrit \emph{ullekha}, 'causing to come forth clearly' (Monier-Williams), before \bn{করা}: to bring a name or a fact out into what you say.
\end{cousinweb}

\begin{grammarlens}[title={What you've built}]
One more word for telling what happened.
\end{grammarlens}

\subsection*{Your turn}
\begin{itemize}
\raggedright
  \item \emph{Say it:} \emph{ullekh kôrā}
  \item \emph{Say it:} \emph{ullekh kôrā}, once more
  \item \emph{Say it:} \emph{prôtibedône sāngbādik sākkhīr nām ullekh kôrlen nā.} — In the report the journalist did not mention the witness's name.
  \item \emph{Recall:} say the Bengali for to announce, then the Bengali for a public statement, then say \emph{ullekh kôrā} again
\end{itemize}

\begin{tcolorbox}[breakable,colback=teal!4,colframe=teal!35!black,title={Before you move on}]
What does \bn{উল্লেখ} \bn{করা} mean? (\textquotedblleft{}To mention\textquotedblright{}.) Say it once more.
\end{tcolorbox}
